\documentclass[11pt,a4paper]{article}

\usepackage[T1]{fontenc}
\usepackage[utf8]{inputenc}
\usepackage{lmodern}
\usepackage{microtype}
\usepackage[a4paper,margin=27mm]{geometry}
\usepackage{amsmath,amssymb,amsthm,mathtools}
\usepackage{booktabs,tabularx,array}
\usepackage{enumitem}
\usepackage[dvipsnames]{xcolor}
\usepackage{parskip}
\usepackage[
  colorlinks=true,
  linkcolor=MidnightBlue,
  urlcolor=MidnightBlue,
  citecolor=MidnightBlue
]{hyperref}
\usepackage{aliascnt}

\setcounter{tocdepth}{1}
\setlist[itemize]{leftmargin=1.5em,itemsep=0.25em,topsep=0.35em}
\setlist[enumerate]{leftmargin=1.7em,itemsep=0.35em,topsep=0.35em}
\emergencystretch=1.5em

\newtheorem{theorem}{Theorem}[section]
\newaliascnt{lemma}{theorem}
\newtheorem{lemma}[lemma]{Lemma}
\aliascntresetthe{lemma}
\theoremstyle{definition}
\newaliascnt{example}{theorem}
\newtheorem{example}[example]{Example}
\aliascntresetthe{example}

\usepackage[capitalize,noabbrev]{cleveref}

\newcommand{\REU}{\operatorname{REU}}
\newcommand{\EU}{\operatorname{EU}}
\newcommand{\V}{\mathcal V}
\newcommand{\D}{\mathcal D}
\newcommand{\U}{\mathcal U}
\newcommand{\dd}{\,\mathrm d}
\newcommand{\ind}{\mathbf 1}
\newcommand{\zero}{\mathbf 0}
\newcommand{\one}{\mathbf 1}

\newenvironment{keyidea}
  {\begin{quote}\color{MidnightBlue}\textbf{Key idea.}\ }
  {\end{quote}}

\newenvironment{teachingnote}
  {\begin{quote}\color{BrickRed}\textbf{Teaching note.}\ }
  {\end{quote}}

\newenvironment{learninggoals}
  {\begin{quote}\color{MidnightBlue}\textbf{Audience and learning goals.}\ }
  {\end{quote}}

\title{When Pointwise REU Choice Undermines Higher-Order REU\\
\large A teaching guide to compression, shadow utility, and small switches}
\author{}
\date{}

\begin{document}
\maketitle

\begin{abstract}
Suppose that every individual decision problem is solved by maximising
risk-weighted expected utility (REU). Must the resulting strategy also
maximise REU after the problems are combined into one reduced lottery over
final utility? In general, no. This note gives a self-contained proof under a
positive-density assumption. A distortion of an attainable event probability
first produces an REU--EU ranking reversal. The two acts are then compressed
into a small interior utility interval. On that interval, the derivative of
the aggregate REU functional is represented by a local shadow utility, which
is nearly affine; its ranking therefore follows the EU difference. A reusable
small-switch lemma turns this strict local disagreement into an exact
higher-order improvement: the first-order gain is of order $\delta$, while the
nonlinear remainder is of order $\delta^2$. The same lemma also supports a
restricted take-or-leave extension, where the shadow utility may have a kink
at zero and three payoff levels replace the compression argument.
\end{abstract}

\begin{learninggoals}
The intended reader is comfortable with finite probability, elementary
measure theory, and Taylor's theorem. The main goals are to distinguish
pointwise from aggregate evaluation, understand why two different small
parameters are needed, and see how one perturbation lemma supports both the
full-domain and take-or-leave results.
\end{learninggoals}

\tableofcontents

\section{The problem and the evaluation convention}

There are two levels of evaluation.

At the \emph{individual level}, a decision problem is presented and REU is
used to choose an act. At the \emph{higher level}, a decision problem is drawn
at random, the strategy chooses an act, and a state is then realised. The
strategy therefore induces a probability distribution over final utility
outcomes.

Throughout the note, the two-stage experiment is \emph{reduced} to that final
utility distribution, and the same risk function is applied to the reduced
lottery. Equivalently, the decision problem and the state are drawn
independently according to the product measure specified below. This is not a
recursive evaluation of an unresolved compound lottery.

The question is:

\begin{quote}
If a strategy chooses an REU-best act in every individual problem, must it
also maximise REU for the reduced lottery generated by the whole strategy?
\end{quote}

In general, no.

\begin{keyidea}
A distorted event probability can make individual REU prefer an act with
lower ordinary expected utility. Compressing that disagreement into a small
utility interval makes the aggregate marginal weight almost constant across
the comparison. The aggregate first-order ranking then follows EU and favours
the act rejected by individual REU.
\end{keyidea}

\begin{example}[A numerical REU--EU reversal]\label{ex:numerical}
Suppose utilities lie in $[0,1]$, an event $A$ has probability $p(A)=1/4$,
and $r(q)=\sqrt q$. Let $g$ pay $1$ on $A$ and $0$ otherwise. Then
\[
 \EU_p(g)=\frac14,
 \qquad
 \REU_{r,p}(g)=r\!\left(\frac14\right)=\frac12.
\]
The constant act $k=3/8$ has both EU and REU equal to $3/8$. Hence
\[
 \REU_{r,p}(g)>\REU_{r,p}(k)
 \qquad\text{but}\qquad
 \EU_p(k)>\EU_p(g).
\]
The main proof constructs this type of reversal for any attainable event
probability $q$ satisfying $r(q)\neq q$.
\end{example}

\section{Formal set-up}

\subsection{Acts, tails, EU, and REU}

Let $\Omega$ be a finite state space, and let $p$ be a full-support
probability measure on $\Omega$. Utilities lie in $[l,h]$, where $l<h$. Write
\[
 \U:=[l,h]^\Omega.
\]
An act is a utility profile $u\in\U$. Throughout the note, $u,v$ denote
generic acts; $A,B,E$ and indexed variants such as $B_i$ denote events in
$\Omega$; $D$ denotes a binary decision problem; $N$ denotes a neighbourhood
of problems or offers; and $S$ denotes a set on which a strategy is switched.
The take-or-leave extension retains these same conventions.

For an act $u$ and a threshold $x$, define its tail probability by
\[
 H_u(x)
 :=p\bigl(\{\omega:u(\omega)>x\}\bigr)
 =\sum_{\omega\in\Omega}p(\omega)
   \ind_{\{u(\omega)>x\}}.
\]
Expected utility has the tail representation
\begin{equation}\label{eq:eu-tail}
 \EU_p(u)
 =l+\int_l^h H_u(x)\dd x.
\end{equation}

Let $r:[0,1]\to[0,1]$ be strictly increasing, with
\[
 r(0)=0,
 \qquad
 r(1)=1.
\]
Risk-weighted expected utility is
\begin{equation}\label{eq:reu-act}
 \REU_{r,p}(u)
 :=l+\int_l^h r\bigl(H_u(x)\bigr)\dd x.
\end{equation}
For decision-theoretic background on REU, see \cite{Buchak2013}. For
background on non-additive and Choquet integration, see
\cite{Denneberg1994,Schmeidler1989}. For the main results, assume
\begin{equation}\label{eq:r-regularity}
 r\in C^2((0,1))
 \qquad\text{and}\qquad
 r'(q)>0\quad\text{for every }q\in(0,1).
\end{equation}

Both EU and REU are positively affine-equivariant. If $\alpha>0$ and the
transformed acts remain in $[l,h]^\Omega$, then
\begin{align*}
 \EU_p(\alpha u+\beta)-\EU_p(\alpha v+\beta)
 &=\alpha\bigl(\EU_p(u)-\EU_p(v)\bigr),\\
 \REU_{r,p}(\alpha u+\beta)-\REU_{r,p}(\alpha v+\beta)
 &=\alpha\bigl(\REU_{r,p}(u)-\REU_{r,p}(v)\bigr).
\end{align*}
Thus a common positive affine transformation preserves both rankings. In the
finite-dimensional act space, both functionals are continuous.

\subsection{Binary decision problems and strategies}

A binary decision problem is an ordered pair
\[
 D=(u,v)\in
 \D:=\U^2.
\]
The order supplies coordinates for the distribution over problems; the
available set is still $\{u,v\}$.

Let $\mu$ be a probability measure on $\D$. For the main theorem, $\mu$ is
absolutely continuous with respect to Lebesgue measure and has density
positive almost everywhere on the interior of $\D$.

A measurable strategy $s$ selects one available act:
\[
 s(u,v)\in\{u,v\}.
\]
It is \emph{pointwise REU-maximising} if
\[
 s(u,v)\in
 \arg\max_{u'\in\{u,v\}}\REU_{r,p}(u')
\]
for every decision problem.

The higher-order experiment is:
\begin{enumerate}
 \item draw $D$ according to $\mu$;
 \item choose $s(D)$;
 \item draw $\omega$ according to $p$;
 \item receive utility $s(D)(\omega)$.
\end{enumerate}
Equivalently, on the product space $(\D\times\Omega,\mu\times p)$ define
\[
 Y_s(D,\omega):=s(D)(\omega).
\]
Its aggregate tail is
\begin{equation}\label{eq:aggregate-tail}
 Q_s(x)
 :=(\mu\times p)(Y_s>x)
 =\int_\D H_{s(D)}(x)\,\mu(\dd D),
\end{equation}
and its higher-order REU is
\begin{equation}\label{eq:higher-value}
 \V(s)
 :=l+\int_l^h r\bigl(Q_s(x)\bigr)\dd x.
\end{equation}

\begin{center}
\renewcommand{\arraystretch}{1.25}
\begin{tabularx}{\textwidth}{>{\bfseries}p{0.22\textwidth}XX}
\toprule
 & Individual act & Whole strategy \\
\midrule
What is random? & The state $\omega$ & First $D$, then $\omega$ \\
Tail probability & $H_u(x)$ & $Q_s(x)=\int H_{s(D)}(x)\,\mu(\dd D)$ \\
Value & $l+\int r(H_u(x))\dd x$ & $l+\int r(Q_s(x))\dd x$ \\
Choice object & One act & A rule over all problems \\
\bottomrule
\end{tabularx}
\end{center}

\subsection{Regularity of the aggregate tail}

\begin{lemma}[Aggregate-tail regularity]\label{lem:aggregate-regularity}
For every measurable strategy $s$ and every $x\in(l,h)$,
\[
 0<Q_s(x)<1.
\]
Moreover, $Q_s$ is continuous on $(l,h)$.
\end{lemma}

\begin{proof}
Let $Y=s(D)(\omega)$ be the final utility under $\mu\times p$. For every
$z\in[l,h]$,
\begin{align*}
 \Pr(Y=z)
 &=\sum_{\omega\in\Omega}p(\omega)
   \mu\bigl(\{D:s(D)(\omega)=z\}\bigr)\\
 &\leq
 \sum_{\omega\in\Omega}p(\omega)
 \left[
   \mu\bigl(\{D:u(\omega)=z\}\bigr)
  +\mu\bigl(\{D:v(\omega)=z\}\bigr)
 \right]
 =0.
\end{align*}
Each set in the last line is a lower-dimensional coordinate slice and has
Lebesgue, hence $\mu$, measure zero. Thus the distribution of $Y$ has no
atoms, and its tail $Q_s(x)=\Pr(Y>x)$ is continuous.

For an interior threshold $x$, there is a non-empty open set of problems on
which both acts have every coordinate above $x$, and another non-empty open
set on which both acts have every coordinate below $x$. Positive density on
the interior gives both sets positive $\mu$-measure. Therefore
$0<Q_s(x)<1$.
\end{proof}

\section{Shadow utility and the small-switch lemma}

The core argument is local. This avoids imposing unnecessary restrictions on
$r'$ near $0$ or $1$.

\subsection{A local shadow utility}

Fix a strategy $s$ and a compact interval
\[
 I=[c,d]\Subset(l,h).
\]
By \cref{lem:aggregate-regularity}, the range of $Q_s$ on $I$ is contained in
a compact subinterval of $(0,1)$. Choose any reference point $z_I\in I$ and
define
\begin{equation}\label{eq:local-shadow-utility}
 W_{s,I}(y)
 :=\int_{z_I}^{y}r'\bigl(Q_s(x)\bigr)\dd x,
 \qquad y\in I.
\end{equation}
The reference point changes $W_{s,I}$ only by an additive constant and hence
does not affect comparisons. The derivative is
\[
 W_{s,I}'(y)=r'\bigl(Q_s(y)\bigr)>0.
\]
For an act $u$ whose coordinates lie in $I$, define its local shadow value by
\begin{equation}\label{eq:local-shadow-value}
 L_{s,I}(u)
 :=\sum_{\omega\in\Omega}p(\omega)W_{s,I}(u(\omega)).
\end{equation}

\begin{lemma}[Shadow-tail identity]\label{lem:shadow-tail}
If the coordinates of $u$ and $v$ lie in $I$, then
\begin{equation}\label{eq:shadow-tail-identity}
 L_{s,I}(v)-L_{s,I}(u)
 =\int_I r'\bigl(Q_s(x)\bigr)
   \bigl(H_v(x)-H_u(x)\bigr)\dd x.
\end{equation}
\end{lemma}

\begin{proof}
For each state,
\[
 W_{s,I}(v(\omega))-W_{s,I}(u(\omega))
 =\int_{u(\omega)}^{v(\omega)}r'\bigl(Q_s(x)\bigr)\dd x.
\]
Summing with weights $p(\omega)$ and interchanging the finite sum and the
integral gives \eqref{eq:shadow-tail-identity}.
\end{proof}

The identity is the first derivative of the aggregate objective written as an
expected-utility comparison. At threshold $x$, a switch changes the chance of
exceeding $x$ by $H_v(x)-H_u(x)$, and the aggregate objective assigns the
marginal weight $r'(Q_s(x))$.

\subsection{One reusable perturbation lemma}

The next result is stated for a generic parameter space so that it can be used
both for binary decision problems and for the take-or-leave extension.

Let $(\Theta,\mathcal F,\mu)$ be an atomless probability space. A strategy
$s$ assigns an act $s(\theta)$ to each parameter $\theta$, and a measurable
map $t$ assigns a feasible alternative act $t(\theta)$. Define
\[
 Q_s(x)=\int_\Theta H_{s(\theta)}(x)\,\mu(\dd\theta)
\]
and define $\V(s)$ by \eqref{eq:higher-value}. Whenever $Q_s$ is
bounded away from $0$ and $1$ on a compact interval $I$, define
$W_{s,I}$ and $L_{s,I}$ by \eqref{eq:local-shadow-utility} and
\eqref{eq:local-shadow-value}.

\begin{lemma}[Small-switch improvement]\label{lem:small-switch}
Let $I=[c,d]\Subset(l,h)$. Suppose there is a measurable set
$N\subseteq\Theta$ with $\mu(N)>0$ and a number $\eta>0$ such that:
\begin{enumerate}[label=(\roman*)]
 \item every coordinate of $s(\theta)$ and $t(\theta)$ lies in $I$ for
       $\theta\in N$;
 \item $Q_s(x)$ is bounded away from $0$ and $1$ on $I$;
 \item
 \[
  L_{s,I}(t(\theta))-L_{s,I}(s(\theta))\geq\eta
  \qquad\text{for every }\theta\in N.
 \]
\end{enumerate}
Then, for every sufficiently small $\delta>0$, there is a measurable
$S\subseteq N$ with $\mu(S)=\delta$ such that the switched strategy
\[
 \widetilde s(\theta)
 =\begin{cases}
   t(\theta),&\theta\in S,\\
   s(\theta),&\theta\notin S
  \end{cases}
\]
satisfies
\[
 \V(\widetilde s)>\V(s).
\]
\end{lemma}

\begin{proof}
Atomlessness gives a measurable $S\subseteq N$ of any sufficiently small
prescribed measure $\delta>0$. Write
\[
 \Delta Q(x):=Q_{\widetilde s}(x)-Q_s(x)
 =\int_S\bigl(H_{t(\theta)}(x)-H_{s(\theta)}(x)\bigr)
   \,\mu(\dd\theta).
\]
Because both changed acts take values in $I$,
\[
 \Delta Q(x)=0\quad\text{for }x\notin I,
 \qquad
 |\Delta Q(x)|\leq\delta\quad\text{for every }x.
\]
For $y,z\in(0,1)$, set
\[
 R_r(y,z):=r(y)-r(z)-r'(z)(y-z).
\]
The exact change in higher-order value is
\begin{align}
 \V(\widetilde s)-\V(s)
 ={}&\int_I r'\bigl(Q_s(x)\bigr)\Delta Q(x)\dd x
 \nonumber\\
 &+\int_I R_r\bigl(Q_s(x)+\Delta Q(x),Q_s(x)\bigr)\dd x.
 \label{eq:exact-local-change}
\end{align}
By Fubini's theorem and \cref{lem:shadow-tail}, the linear term equals
\begin{align*}
 \int_S
 \bigl[L_{s,I}(t(\theta))-L_{s,I}(s(\theta))\bigr]
 \,\mu(\dd\theta)
 \geq\eta\delta.
\end{align*}

Choose $\kappa>0$ such that
\[
 \kappa\leq Q_s(x)\leq1-\kappa
 \qquad(x\in I).
\]
For $\delta<\kappa/2$, both arguments of $R_r$ and the segment between them
lie in $[\kappa/2,1-\kappa/2]$. If
\[
 M:=\sup_{q\in[\kappa/2,1-\kappa/2]}|r''(q)|,
\]
Taylor's theorem gives
\[
 \left|R_r\bigl(Q_s(x)+\Delta Q(x),Q_s(x)\bigr)\right|
 \leq\frac{M}{2}\delta^2.
\]
Consequently,
\[
 \V(\widetilde s)-\V(s)
 \geq\eta\delta-\frac{M}{2}|I|\delta^2>0
\]
for every sufficiently small $\delta>0$.
\end{proof}

\begin{keyidea}
Once a strict shadow-value improvement exists on a positive-measure set, the
rest is automatic. The average first-order gain is proportional to the size
$\delta$ of the switched set, while the nonlinear error is proportional to
$\delta^2$.
\end{keyidea}

\section{The full-domain theorem}

\begin{theorem}[Positive-density self-undermining]\label{thm:main}
Suppose there is a non-empty proper event $A\subsetneq\Omega$ such that
\[
 r\bigl(p(A)\bigr)\neq p(A).
\]
Let $\mu$ be absolutely continuous with density positive almost everywhere on
the interior of $\D$. Then every measurable pointwise REU-maximising strategy
$s$ is strictly suboptimal according to higher-order REU.

More precisely, there are a measurable set $S\subseteq\D$ with $\mu(S)>0$
and a strategy $\widetilde s$ such that:
\begin{enumerate}[label=(\roman*)]
 \item $\widetilde s=s$ outside $S$;
 \item on $S$, $\widetilde s$ chooses the individually REU-inferior act;
 \item $\V(\widetilde s)>\V(s)$.
\end{enumerate}
\end{theorem}

\subsection{An attainable event distortion creates a reversal}

\begin{lemma}[REU--EU reversal]\label{lem:reu-eu-reversal}
If $A$ is a non-empty proper event and
$r(p(A))\neq p(A)$, then there are acts $u^\star,v^\star\in\U$ such that
\begin{equation}\label{eq:reu-eu-reversal}
 \REU_{r,p}(u^\star)>\REU_{r,p}(v^\star)
 \qquad\text{and}\qquad
 \EU_p(v^\star)>\EU_p(u^\star).
\end{equation}
\end{lemma}

\begin{proof}
Let $q=p(A)$ and let $g$ pay $h$ on $A$ and $l$ on $A^c$. Then
\[
 \EU_p(g)=l+(h-l)q,
 \qquad
 \REU_{r,p}(g)=l+(h-l)r(q).
\]
A constant act $k$ has both EU and REU equal to $k$.

If $r(q)>q$, choose
\[
 l+(h-l)q<k<l+(h-l)r(q)
\]
and set $u^\star=g$, $v^\star=k$. If $r(q)<q$, choose
\[
 l+(h-l)r(q)<k<l+(h-l)q
\]
and set $u^\star=k$, $v^\star=g$. In either case
\eqref{eq:reu-eu-reversal} holds.
\end{proof}

\subsection{Compression converts the EU ranking into a shadow ranking}

Fix a pointwise REU-maximising strategy $s$ and an interior utility level
$z\in(l,h)$. Let $u^\star,v^\star$ be the acts supplied by
\cref{lem:reu-eu-reversal}. For $\varepsilon>0$, define
\begin{equation}\label{eq:compressed-acts}
 u_\varepsilon=z\one+\varepsilon(u^\star-z\one),
 \qquad
 v_\varepsilon=z\one+\varepsilon(v^\star-z\one).
\end{equation}

\begin{lemma}[Compression lemma]\label{lem:compression}
There are a compact interval $I\Subset(l,h)$ containing $z$ in its interior
and a number $\bar\varepsilon>0$ such that, for every
$0<\varepsilon<\bar\varepsilon$, all coordinates of
$u_\varepsilon,v_\varepsilon$ lie in $\operatorname{int} I$ and
\begin{equation}\label{eq:compressed-reversal}
 \REU_{r,p}(u_\varepsilon)>\REU_{r,p}(v_\varepsilon)
 \qquad\text{but}\qquad
 L_{s,I}(v_\varepsilon)>L_{s,I}(u_\varepsilon).
\end{equation}
\end{lemma}

\begin{proof}
Choose $\rho>0$ with $I=[z-\rho,z+\rho]\Subset(l,h)$. For sufficiently small
$\varepsilon$, all compressed coordinates lie in $\operatorname{int} I$.
Positive affine equivariance gives
\begin{align*}
 \REU_{r,p}(u_\varepsilon)-\REU_{r,p}(v_\varepsilon)
 &=\varepsilon
   \bigl(\REU_{r,p}(u^\star)-\REU_{r,p}(v^\star)\bigr)>0,\\
 \EU_p(v_\varepsilon)-\EU_p(u_\varepsilon)
 &=\varepsilon
   \bigl(\EU_p(v^\star)-\EU_p(u^\star)\bigr)>0.
\end{align*}
Use $z$ as the reference point in \eqref{eq:local-shadow-utility}. Since
$W_{s,I}$ is differentiable at $z$,
\[
 \frac{W_{s,I}\bigl(z+\varepsilon(y-z)\bigr)-
 W_{s,I}(z)}{\varepsilon(y-z)}\longrightarrow W_{s,I}'(z)
\]
so 
\[
 \frac{W_{s,I}\bigl(z+\varepsilon(y-z)\bigr)-
 W_{s,I}(z)}{\varepsilon(y-z)}\longrightarrow W_{s,I}'(z)
\]
uniformly for $y\in[l,h]$. Summing over the finite state space yields
\begin{align}
 L_{s,I}(v_\varepsilon)-L_{s,I}(u_\varepsilon)
 ={}&\varepsilon W_{s,I}'(z)
 \bigl(\EU_p(v^\star)-\EU_p(u^\star)\bigr)+o(\varepsilon).
 \label{eq:shadow-compression}
\end{align}
By \cref{lem:aggregate-regularity} and \eqref{eq:r-regularity},
\[
 W_{s,I}'(z)=r'\bigl(Q_s(z)\bigr)>0.
\]
The leading coefficient in \eqref{eq:shadow-compression} is therefore
positive, which proves the shadow inequality for all sufficiently small
$\varepsilon$.
\end{proof}

\begin{teachingnote}
Compression does not make REU itself become EU. It makes the aggregate
marginal weight $r'(Q_s(x))$ nearly constant over the thresholds affected by
this particular comparison. That is why the local shadow ranking follows the
EU difference.
\end{teachingnote}

\subsection{Proof of the theorem}

\begin{proof}[Proof of \cref{thm:main}]
Fix a measurable pointwise REU-maximising strategy $s$. Choose
$\varepsilon>0$, $I$, and the compressed pair
$(u_\varepsilon,v_\varepsilon)$ as in \cref{lem:compression}.

Both REU and the local shadow value are continuous in the acts. The
inequalities in \eqref{eq:compressed-reversal} are strict, so there are an
open neighbourhood $N$ of $(u_\varepsilon,v_\varepsilon)$ and a number
$\eta>0$ such that
\[
 \overline N\subset
 ((\operatorname{int} I)^\Omega)^2
 \subseteq\operatorname{int}(\D),
\]
and, for every $(u,v)\in N$,
\begin{equation}\label{eq:uniform-reversal}
 \REU_{r,p}(u)>\REU_{r,p}(v),
 \qquad
 L_{s,I}(v)-L_{s,I}(u)\geq\eta.
\end{equation}
Every pointwise REU-maximising strategy therefore selects $u$ throughout
$N$.

Positive density gives $\mu(N)>0$, and absolute continuity makes $\mu$
atomless. Apply \cref{lem:small-switch} with parameter space $\Theta=\D$,
baseline act $s(u,v)=u$, and alternative act $t(u,v)=v$ on $N$. It produces a
measurable set $S\subseteq N$ of positive measure such that switching from
$u$ to $v$ on $S$ strictly increases higher-order REU. The switched act is
strictly REU-inferior on every problem in $S$, so all three conclusions of the
theorem follow.
\end{proof}

\subsection{Why two small parameters appear}

\begin{center}
\renewcommand{\arraystretch}{1.25}
\begin{tabularx}{\textwidth}{>{\bfseries}p{0.18\textwidth}XX}
\toprule
Parameter & What it makes small & Its role \\
\midrule
$\varepsilon$ & The utility spread of the compressed acts & Makes the local shadow utility nearly affine, so the shadow comparison follows EU \\
$\delta=\mu(S)$ & The probability of the switched set & Makes the nonlinear $O(\delta^2)$ remainder smaller than the guaranteed $O(\delta)$ gain \\
\bottomrule
\end{tabularx}
\end{center}

\section{What the assumptions are doing}

\begin{center}
\renewcommand{\arraystretch}{1.3}
\begin{tabularx}{\textwidth}{>{\raggedright\arraybackslash\bfseries}p{0.27\textwidth}X}
\toprule
Assumption & Role in the proof \\
\midrule
Finite state space & Makes the act space finite-dimensional and lets the shadow expansion be summed state by state. \\
Full support of $p$ & Gives every non-empty proper event a probability strictly between $0$ and $1$. \\
$r'(q)>0$ & Makes the local shadow utility increasing, so a positive EU difference remains positive after compression. \\
$r\in C^2((0,1))$ & Gives a uniform quadratic bound on the nonlinear remainder over the relevant compact range. \\
Absolute continuity of $\mu$ & Makes $\mu$ atomless and rules out atoms in the aggregate utility distribution generated by binary pairs. \\
Positive density on the interior & Ensures that the open neighbourhood of strict reversals has positive probability. \\
$r(p(A))\neq p(A)$ & Produces the initial REU--EU reversal at an event probability attainable under $p$. \\
\bottomrule
\end{tabularx}
\end{center}

The last condition is more precise than saying merely that $r$ is nonlinear.
For a fixed finite state space, every tail probability $H_u(x)$ is the
probability of an event. If
\[
 r\bigl(p(E)\bigr)=p(E)
 \qquad\text{for every event }E\subseteq\Omega,
\]
then $r(H_u(x))=H_u(x)$ at every threshold of every act, so
\[
 \REU_{r,p}(u)=\EU_p(u)
 \qquad\text{for every act }u.
\]
A risk function may therefore be nonlinear away from the finite set of
attainable event probabilities without generating the reversal used here.

The assumptions can also be weakened locally. The proof needs positive mass
only on a neighbourhood of the constructed reversal, and it needs a quadratic
remainder bound only on the compact probability range visited by the switch.
The stated conditions are convenient transparent sufficient conditions.

\section{A take-or-leave extension}

The full-domain theorem uses arbitrary nearby pairs of acts. We now consider a
much thinner support in which every problem offers only the common baseline
$\zero$ and one act $u$. The act space remains $\U$; only the menu changes.

\subsection{Set-up and the kink at zero}

Assume $l<0<h$. For each $u\in\U$, the decision problem is $\{\zero,u\}$,
where $\zero$ is the constant-zero act. Let $\mu$ be an absolutely continuous
probability measure on $\U$ whose density is positive almost everywhere on an
open neighbourhood $N_0$ of $\zero$.

A take-or-leave strategy is a measurable map satisfying
\[
 s(u)\in\{\zero,u\}.
\]
It is pointwise REU-maximising when it accepts $u$ if
$\REU_{r,p}(u)>0$, rejects it if $\REU_{r,p}(u)<0$, and may choose either at a
tie. Its aggregate tail and higher-order value are defined by
\[
 Q_s(x)=\int_\U H_{s(u)}(x)\,\mu(\dd u),
 \qquad
 \V(s)=l+\int_l^h r\bigl(Q_s(x)\bigr)\dd x.
\]

Rejected offers generate the outcome zero, so the final distribution may have
an atom at zero.

\begin{lemma}[The aggregate distribution around zero]\label{lem:zero-mass}
For every pointwise REU-maximising take-or-leave strategy,
\[
 \Pr(Y>0)>0,
 \qquad
 \Pr(Y=0)>0,
 \qquad
 \Pr(Y<0)>0.
\]
The final utility distribution has no atoms away from zero. Consequently, with
\[
 Q_s(0^+):=\lim_{x\downarrow0}Q_s(x)=\Pr(Y>0),
 \qquad
 Q_s(0^-):=\lim_{x\uparrow0}Q_s(x)=\Pr(Y\geq0),
\]
we have
\begin{equation}\label{eq:zero-tail-inequalities}
 0<Q_s(0^+)<Q_s(0^-)<1.
\end{equation}
\end{lemma}

\begin{proof}
Because $N_0$ is a neighbourhood of $\zero$, it contains sufficiently small
constant positive and constant negative acts. An open set of strictly positive
acts is strictly preferred to $\zero$ and is accepted, so positive outcomes
occur with positive probability. An open set of strictly negative acts is
strictly worse than $\zero$ and is rejected, so the outcome zero has positive
probability.

To obtain negative outcomes, choose a non-empty proper event $B$ and let
$q=p(B)$. For small $c,d>0$, consider the act paying $c$ on $B$ and $-d$ on
$B^c$. Its REU is
\[
 r(q)c-\bigl(1-r(q)\bigr)d.
\]
Choose $d/c$ small enough to make this expression positive, and then scale the
act so that it lies in $N_0$. On a sufficiently small open neighbourhood, REU
remains positive and every payoff on $B^c$ remains negative. Those acts are
accepted, so negative outcomes occur with positive probability. Positive
density gives all three open sets positive $\mu$-measure.

For $z\neq0$, the event $Y=z$ is contained in
\[
 \bigcup_{\omega\in\Omega}\{u:u(\omega)=z\},
\]
a finite union of Lebesgue-null coordinate slices. Thus there are no atoms
away from zero. The identities and inequalities in
\eqref{eq:zero-tail-inequalities} now follow from the definitions of the
one-sided tail limits.
\end{proof}

By \cref{lem:zero-mass}, choose $\rho>0$ such that
\[
 I_0=[-\rho,\rho]\Subset(l,h)
\]
and $Q_s(x)$ is bounded away from $0$ and $1$ on $I_0$. Use the same local
shadow notation as before, choosing zero as the reference point:
\begin{equation}\label{eq:zero-shadow}
 W_{s,I_0}(y):=\int_0^y r'\bigl(Q_s(x)\bigr)\dd x,
 \qquad y\in I_0.
\end{equation}
The value assigned to the integrand at the single point $x=0$ is immaterial.
The one-sided slopes are
\begin{equation}\label{eq:one-sided-slopes}
 w_+:=r'\bigl(Q_s(0^+)\bigr)>0,
 \qquad
 w_-:=r'\bigl(Q_s(0^-)\bigr)>0.
\end{equation}
For acts with coordinates in $I_0$, retain the interval-indexed shadow value
\[
 L_{s,I_0}(u):=\sum_{\omega\in\Omega}p(\omega)W_{s,I_0}(u(\omega)).
\]

The directional shadow value at the zero act is the kinked linear functional
\begin{equation}\label{eq:kinked-functional}
 \Lambda_s(u)
 :=w_+\sum_{u(\omega)>0}p(\omega)u(\omega)
  +w_-\sum_{u(\omega)<0}p(\omega)u(\omega).
\end{equation}

\begin{lemma}[Kinked shadow expansion]\label{lem:kinked-expansion}
For every fixed act $u\in\U$,
\begin{equation}\label{eq:kinked-expansion}
 L_{s,I_0}(\varepsilon u)-L_{s,I_0}(\zero)
 =\varepsilon\Lambda_s(u)+o(\varepsilon)
 \qquad(\varepsilon\downarrow0).
\end{equation}
\end{lemma}
\begin{proof}
If $y>0$, the right derivative of $W_{s,I_0}$ at zero is $w_+$, while if
$y<0$, the left derivative is $w_-$. Hence
\[
 \frac{W_{s,I_0}(\varepsilon y)-W_{s,I_0}(0)}{\varepsilon}
 \longrightarrow
 \begin{cases}
  w_+y,&y>0,\\
  w_-y,&y<0,\\
  0,&y=0.
 \end{cases}
\]
Summing over the finite state space gives \eqref{eq:kinked-expansion}.
\end{proof}

\subsection{A geometric sign lemma}

Let
\[
 C_{++-}:=\{\xi\in\mathbb R^3:\xi_1>\xi_2>0>\xi_3\},
\]
\[
 C_{+--}:=\{\xi\in\mathbb R^3:\xi_1>0>\xi_2>\xi_3\}.
\]

\begin{lemma}[Agreement on either sign cone]\label{lem:sign-cone}
Let $C$ be either $C_{++-}$ or $C_{+--}$, and let
$\alpha,\beta\in(0,\infty)^3$. If
\[
 (\alpha\cdot\xi)(\beta\cdot\xi)\geq0
 \qquad\text{for every }\xi\in C,
\]
then $\alpha$ and $\beta$ are positive scalar multiples of one another.
\end{lemma}

\begin{proof}
The plane $\{\xi:\alpha\cdot\xi=0\}$ intersects either cone in a non-empty
relatively open set. For $C_{++-}$, choose $\xi_1>\xi_2>0$ and solve for a
negative $\xi_3$; for $C_{+--}$, choose $0>\xi_2>\xi_3$ and solve for a
positive $\xi_1$.

If $\beta$ were not proportional to $\alpha$, its restriction to that plane
would not vanish identically. Hence there would be a point $\xi^0$ in the
intersection with $\beta\cdot\xi^0\neq0$. A sufficiently small perturbation
of $\xi^0$ remains in the open cone and crosses the plane
$\alpha\cdot\xi=0$ while preserving the sign of $\beta\cdot\xi$. On one side,
the two linear forms would have opposite signs, contradicting the hypothesis.
Since both coefficient vectors are positive, the scalar multiple is positive.
\end{proof}

\subsection{The restricted-support theorem}

\begin{theorem}[Rich take-or-leave problems]\label{thm:take-or-leave}
Let $|\Omega|\geq3$, let $p$ have full support, and suppose $l<0<h$. Let $r$
satisfy \eqref{eq:r-regularity}, with $r(0)=0$ and $r(1)=1$. Suppose that the
distribution $\mu$ over offered acts is absolutely continuous and has density
positive almost everywhere on an open neighbourhood of the zero act.

If
\[
 r\bigl(p(A)\bigr)\neq p(A)
\]
for some non-empty proper event $A\subsetneq\Omega$, then every pointwise
REU-maximising take-or-leave strategy is strictly suboptimal according to
higher-order REU.
\end{theorem}

\begin{proof}
Fix a pointwise REU-maximising strategy $s$ and the local objects constructed
above. Choose a partition $B_1,B_2,B_3$ of $\Omega$ into three non-empty
events as follows. If $A^c$ contains at least two states, set $B_1=A$ and
split $A^c$ into $B_2,B_3$. If $A^c$ is a singleton, split $A$ into non-empty
$B_1,B_2$ and set $B_3=A^c$. Write
\[
 q_i:=p(B_i),
 \qquad
 q_1+q_2+q_3=1.
\]
The construction ensures that at least one of
\[
 r(q_1)\neq q_1,
 \qquad
 r(q_1+q_2)\neq q_1+q_2
\]
holds.

For $\xi_1>\xi_2>\xi_3$, consider the three-level act
\[
 u^\xi=\xi_1\ind_{B_1}+\xi_2\ind_{B_2}+\xi_3\ind_{B_3}.
\]
Its REU can be written first in incremental form and then as a linear form:
\begin{align}
 \REU_{r,p}(u^\xi)
 &={}\xi_3+r(q_1+q_2)(\xi_2-\xi_3)+r(q_1)(\xi_1-\xi_2)\nonumber\\
 &=\alpha_1\xi_1+\alpha_2\xi_2+\alpha_3\xi_3,
 \label{eq:three-level-reu}
\end{align}
where
\begin{equation}\label{eq:three-reu-weights}
 \alpha_1=r(q_1),
 \qquad
 \alpha_2=r(q_1+q_2)-r(q_1),
 \qquad
 \alpha_3=1-r(q_1+q_2).
\end{equation}
All three coefficients are positive and sum to one.

On the cone $C_{++-}$, the directional shadow value is
\[
 \Lambda_s(u^\xi)
 =w_+q_1\xi_1+w_+q_2\xi_2+w_-q_3\xi_3.
\]
If REU and the directional shadow value never had opposite signs there,
\cref{lem:sign-cone} would imply
\[
 (\alpha_1,\alpha_2,\alpha_3)
 \ \parallel\
 (w_+q_1,w_+q_2,w_-q_3),
\]
and hence
\begin{equation}\label{eq:first-ratio}
 \frac{\alpha_1}{q_1}=\frac{\alpha_2}{q_2}.
\end{equation}
On the cone $C_{+--}$, the same argument would imply
\[
 (\alpha_1,\alpha_2,\alpha_3)
 \ \parallel\
 (w_+q_1,w_-q_2,w_-q_3),
\]
and hence
\begin{equation}\label{eq:second-ratio}
 \frac{\alpha_2}{q_2}=\frac{\alpha_3}{q_3}.
\end{equation}

If no sign reversal existed in either cone, then
\[
 \frac{\alpha_1}{q_1}
 =\frac{\alpha_2}{q_2}
 =\frac{\alpha_3}{q_3}
 =:c.
\]
Since both the $\alpha_i$ and the $q_i$ sum to one, $c=1$. Therefore
$\alpha_i=q_i$ for every $i$, which implies
\[
 r(q_1)=q_1,
 \qquad
 r(q_1+q_2)=q_1+q_2.
\]
This contradicts the way the partition was chosen. Hence there is an
$\xi\in C_{++-}\cup C_{+--}$ such that
\begin{equation}\label{eq:take-leave-sign-reversal}
 \REU_{r,p}(u^\xi)\,\Lambda_s(u^\xi)<0.
\end{equation}

Positive homogeneity gives
\[
 \REU_{r,p}(\varepsilon u^\xi)
 =\varepsilon\REU_{r,p}(u^\xi).
\]
By \cref{lem:kinked-expansion}, for all sufficiently small
$\varepsilon>0$, the exact shadow difference
\[
 L_{s,I_0}(\varepsilon u^\xi)-L_{s,I_0}(\zero)
\]
has the sign of $\Lambda_s(u^\xi)$. Choose $\varepsilon$ small enough that the
offered act $\bar u:=\varepsilon u^\xi$ lies in $N_0$ and all its coordinates lie
in $I_0$.

If REU is positive and the shadow difference is negative, the pointwise
strategy accepts $\bar u$ while the shadow value favours rejection. If REU is
negative and the shadow difference is positive, the pointwise strategy rejects
$\bar u$ while the shadow value favours acceptance. Because both inequalities
are strict and both functionals are continuous, there are an open
neighbourhood $N\subseteq N_0$ of $\bar u$ and a number $\eta>0$ on which the
same two rankings persist.

Define the alternative to be the other member of each take-or-leave menu: in
the acceptance case set $t(u)=\zero$ for $u\in N$, and in the rejection case
set $t(u)=u$ for $u\in N$. Then
\[
 L_{s,I_0}(t(u))-L_{s,I_0}(s(u))\geq\eta
 \qquad(u\in N),
\]
while $t(u)$ is strictly REU-inferior pointwise. Shrink $N$ so that all
relevant payoffs lie in $I_0$. Positive density gives $\mu(N)>0$, and absolute
continuity gives atomlessness. Applying \cref{lem:small-switch} with parameter
space $\Theta=\U$ produces a positive-measure switch set $S\subseteq N$ that
strictly increases higher-order REU.
\end{proof}

\begin{keyidea}
The full-domain theorem and the take-or-leave theorem differ only in how they
find a shadow-value reversal. In the full domain, compression makes the local
shadow comparison follow EU. In the take-or-leave domain, zero can create a
kink, and three payoff levels show that the two one-sided shadow weights cannot
reproduce distorted REU decision weights in both sign patterns.
\end{keyidea}

\section{Comparison, checks, and conclusion}

\begin{center}
\renewcommand{\arraystretch}{1.25}
\begin{tabularx}{\textwidth}{>{\raggedright\arraybackslash\bfseries}p{0.24\textwidth}XX}
\toprule
 & Full binary domain & Take-or-leave domain \\
\midrule
Available comparison & Arbitrary nearby pair $\{u,v\}$ & Only $\{\zero,u\}$ \\
Compression point & Any interior $z$ & The fixed baseline $\zero$ \\
Local shadow shape & Smooth and nearly affine & Potentially kinked at zero \\
Reversal device & REU--EU reversal plus compression & Three-level coefficient comparison \\
Common final step & \multicolumn{2}{l}{Small-switch lemma: an $O(\delta)$ gain dominates an $O(\delta^2)$ error} \\
\bottomrule
\end{tabularx}
\end{center}

The proof architecture is therefore:
\[
 \boxed{
 \begin{gathered}
  \text{find a strict REU--shadow reversal}\\[2pt]
  \Downarrow\\[-1pt]
  \text{switch on a small positive-measure set}\\[2pt]
  \Downarrow\\[-1pt]
  \text{higher-order REU rises}
 \end{gathered}}
\]

\subsection*{Reader checks}

\begin{enumerate}
 \item Verify \eqref{eq:eu-tail} directly by writing
       $u(\omega)=l+\int_l^h\ind_{\{u(\omega)>x\}}\dd x$ and summing over
       states.
 \item Explain why the hypothesis $r(p(A))\neq p(A)$ concerns an attainable
       event probability rather than global nonlinearity of $r$.
 \item In \cref{lem:small-switch}, identify exactly where the orders
       $O(\delta)$ and $O(\delta^2)$ enter and why neither can be replaced by
       the compression parameter $\varepsilon$.
 \item Locate the step in the take-or-leave proof that requires the three
       non-empty events $B_1,B_2,B_3$. What coefficient comparison is
       unavailable when $|\Omega|=2$?
\end{enumerate}

The central lesson is not simply that nonlinear weighting creates a
preference reversal. It is that a rich distribution of decision problems
lets that reversal be placed where the aggregate objective has a different
local geometry. The shadow utility exposes that geometry, and the small-switch
lemma converts the local disagreement into an exact global improvement.

\begin{thebibliography}{9}
\bibitem{Buchak2013}
Lara Buchak.
\newblock \emph{Risk and Rationality}.
\newblock Oxford University Press, 2013.

\bibitem{Denneberg1994}
Dieter Denneberg.
\newblock \emph{Non-Additive Measure and Integral}.
\newblock Kluwer Academic Publishers, 1994.

\bibitem{Schmeidler1989}
David Schmeidler.
\newblock Subjective probability and expected utility without additivity.
\newblock \emph{Econometrica}, 57(3):571--587, 1989.
\end{thebibliography}

\end{document}
